% ---------------------------------------------------------------------------
% Shared preamble: fonts, geometry, colour convention, pgfplots styles.
% The colour convention is deliberate: the same quantity always gets the same
% colour throughout the report (see Section "Conventions" of Chapter 2).
% ---------------------------------------------------------------------------
\usepackage[T1]{fontenc}
\usepackage[utf8]{inputenc}
\usepackage{newtxtext,newtxmath}
\usepackage[english]{babel}
\usepackage[a4paper,top=2.6cm,bottom=2.6cm,left=2.7cm,right=2.7cm]{geometry}
\usepackage{setspace}
\setstretch{1.17}

\usepackage{booktabs}
\usepackage{longtable}
\usepackage{multirow}
\usepackage{array}
\usepackage{tabularx}
\usepackage{siunitx}
\sisetup{group-separator={\,},detect-all}
\usepackage{graphicx}
\usepackage{xcolor}
\usepackage{tikz}
\usepackage{pgfplots}
\pgfplotsset{compat=1.18}
\usepgfplotslibrary{groupplots}
\usetikzlibrary{arrows.meta,positioning,fit,calc,decorations.pathreplacing,backgrounds,shapes.geometric}
\usepackage{caption}
\captionsetup{font=small,labelfont=bf,width=0.94\textwidth}
\usepackage{enumitem}
\setlist{itemsep=2pt,topsep=3pt,parsep=0pt}
\usepackage[hidelinks]{hyperref}
\usepackage{titlesec}
\usepackage{fancyhdr}
\setlength{\headheight}{14pt}
\usepackage{emptypage}

% --- Headers ---------------------------------------------------------------
\pagestyle{fancy}
\fancyhf{}
\fancyhead[L]{\small\nouppercase{\leftmark}}
\fancyhead[R]{\small\thepage}
\renewcommand{\headrulewidth}{0.4pt}
\fancypagestyle{plain}{\fancyhf{}\fancyhead[R]{\small\thepage}\renewcommand{\headrulewidth}{0pt}}

% --- Chapter/section titles ------------------------------------------------
\titleformat{\chapter}[display]
  {\normalfont\Large\bfseries}{\normalsize\MakeUppercase{\chaptertitlename}\ \thechapter}{8pt}{\LARGE}
\titlespacing*{\chapter}{0pt}{-10pt}{24pt}

% --- THE COLOUR CONVENTION -------------------------------------------------
% Never redefine these locally. The reader is meant to learn them once.
\definecolor{clbase}{RGB}{108,125,145}   % persistence baseline
\definecolor{clmodel}{RGB}{31,56,100}    % the fine-tuned model
\definecolor{clgain}{RGB}{178,84,32}     % a positive gap / a gain
\definecolor{clloss}{RGB}{150,60,60}     % a negative gap
\definecolor{clnonsig}{RGB}{186,192,199} % measured but not significant
\definecolor{clink}{RGB}{70,150,170}     % "same site, other cell"
\definecolor{clmove}{RGB}{216,132,32}    % "real site change"
\definecolor{clrule}{RGB}{140,148,158}
\definecolor{clsoft}{RGB}{240,242,245}
% The four behavioural groups, ordered from sedentary to mobile.
\definecolor{clGzero}{RGB}{70,120,160}
\definecolor{clGone}{RGB}{90,150,80}
\definecolor{clGtwo}{RGB}{214,148,44}
\definecolor{clGthree}{RGB}{170,60,50}

% --- Shared pgfplots style -------------------------------------------------
\pgfplotsset{
  rep/.style={
    width=\linewidth, height=5.4cm,
    font=\small,
    axis line style={clrule},
    tick align=outside, tick style={clrule},
    grid=major, grid style={clnonsig!45,line width=0.3pt},
    every axis label/.append style={font=\small},
    legend style={font=\small,draw=clrule,fill=white,fill opacity=0.9,text opacity=1},
    label style={font=\small}, tick label style={font=\small},
  },
  repbar/.style={rep, ybar, /pgf/bar width=9pt, ymajorgrids, xmajorgrids=false,
                 enlarge x limits=0.12, nodes near coords,
                 /pgf/number format/.cd, fixed, fixed zerofill, precision=1, /pgfplots/.cd,
                 every node near coord/.append style={font=\tiny,color=black,rotate=90,anchor=west}},
  rephbar/.style={rep, xbar, /pgf/bar width=8pt, xmajorgrids, ymajorgrids=false},
}
% A manual legend row, placed ABOVE a tikzpicture so it can never collide with
% a plot title. \swatch{colour} draws the key; \replegend collects a row.
% Shrink a box to the text width only if it is wider; never enlarge it.
\newcommand{\fitwidth}[1]{\resizebox{\ifdim\width>\textwidth \textwidth\else \width\fi}{!}{#1}}
\newcommand{\swatch}[1]{\tikz[baseline=-0.55ex]{\fill[#1] (0,0) rectangle (0.30,0.17);}}
\newcommand{\replegend}[1]{\par\smallskip{\footnotesize #1\par}\smallskip}

% --- Small semantic macros used in the prose -------------------------------
\newcommand{\pt}{\,pt}
\newcommand{\cid}[1]{\texttt{#1}}
% Datasets are referred to by a descriptive name, never a code: the reader
% should never have to look up what "D3" was.
\newcommand{\Draw}{the \textsc{raw day}}
\newcommand{\Drawbig}{the \textsc{raw day} at 5{,}000 users}
\newcommand{\Dfilt}{the \textsc{filtered day}}
\newcommand{\Dmob}{the \textsc{mobile-user day}}
\newcommand{\Dlev}{the \textsc{levelled corpus}}
\newcommand{\Dtwok}{the \textsc{2{,}000-user day}}
\newcommand{\Deleven}{the \textsc{eleven-day set}}
% Short forms for axis labels and table cells, where "the" would be wrong.
\newcommand{\sraw}{\textsc{raw day}}
\newcommand{\srawbig}{\textsc{raw day}, 5{,}000 users}
\newcommand{\sfilt}{\textsc{filtered day}}
\newcommand{\smob}{\textsc{mobile-user day}}
\newcommand{\slev}{\textsc{levelled corpus}}
\newcommand{\stwok}{\textsc{2{,}000-user day}}
\newcommand{\seleven}{\textsc{eleven-day set}}
\newcommand{\supervisor}{Robert Best\'ak}
\newcommand{\host}{Czech Technical University in Prague (\v{C}VUT)}
